\clearpage
\begingroup
\revj
\section{The free massless field on Minkowski space}\label{sec:minkowski}

\noindent\motiv{Chapter~\ref{sec:green} reduces the axioms and covariance to Green data. This
chapter constructs Green data for the massless wave operator $P=\partial_t^2-\nabla^2$ on $3+1$
Minkowski space, the operator $P=-\square+m^2$ of Remark~\ref{rem:dictionary} with $m=0$. The
retarded and advanced Green operators are given by the \revm{retarded potential and its advanced
counterpart}, and the cutoffs are slabs
between parallel spacelike hyperplanes.}

\medskip\noindent\motiv{For the free massless field the hypothesis of the introduction, the existence of Green
operators, thereby becomes a theorem. The CCR net is generating and causal, and it has the
time-slice property for these slabs, which \revk{implies the time-slice axiom in the form
of~\cite{dimock} (Example~\ref{ex:gh-index})} for flat Cauchy surfaces
(Remark~\ref{rem:mink-slabs}\revk{, on paper}). Its algebra is primitive, given an irreducible representation, and
it is covariant under the orthochronous Poincar\'e group.}

\medskip\noindent\motiv{Remark~\ref{rem:green-continuum} sketches the general case with all Cauchy temporal functions
as cutoffs; that larger family is not formalized. \revm{Chapter~\ref{sec:mass} adds the term $m^2$, for every real $m^2$.}}

\subsection{Minkowski space and the wave operator}

\begin{definition}[Minkowski space]\label{def:mink}
Points are $p=(p^0,\vec p)\in M=\R\times\R^3$, with the Euclidean norm $|\cdot|$ on $\R^3$ and the
Minkowski form $\eta(u,w)=-u^0w^0+\vec u\cdot\vec w$ (mostly plus); $e_0=(1,\vec0)$. A vector $v$ is
\emph{future causal} if $v^0\geq|\vec v|$, and $p\leq q$ iff $q-p$ is future causal;
$I^+=\{v:v^0>|\vec v|\}$ and $I^-=-I^+$ are the open future and past cones.

$\cD=C_c^\infty(M)$ and $\cF=C^\infty(M)$ are the real test functions and smooth functions, $\iota$ is
the inclusion, and $\supp h=\overline{\{h\neq0\}}$ is the closed support. The wave operator
$P=\partial_0^2-\sum_{i=1}^3\partial_i^2=-\eta^{\mu\nu}\partial_\mu\partial_\nu$ acts on $\cF$ and on $\cD$,
and $\langle f,h\rangle=\int_Mf(p)h(p)\,dp$ with Lebesgue measure on $M$. \revm{In coordinates
$(t,x)$ the spatial variable $x\in\R^3$ carries no arrow.}
\end{definition}

\begin{remark}[Closed supports]\label{rem:mink-supports}
With the closed support, $\cD(O)$ is $C_c^\infty(O)$ for open $O$, as in~\cite{dimock}: a test
function localized in $O$ vanishes near the boundary of $O$.

With the open support $\{h\neq0\}$ instead, a test function could be nonzero up to the boundary, and
two things would change, as the double cone $D=\{|t|+|\vec x|<1\}$ shows. A test function positive
exactly on $D\cap\{t>0\}$ would give a Weyl element outside the usual algebra of that region. And
$D$ would no longer depend on a thin layer $D\cap\{|t|<\ell\}$ around its base, since one test
function could fill all of $D$.

\revm{\artefact{The formalization uses the closed support.}}
\end{remark}

\begin{proposition}[The Minkowski data]\label{prop:mink-basic}
The objects of Definition~\ref{def:mink} satisfy (G1)--(G5) of Definition~\ref{def:green-data}. In
detail:
\begin{enumerate}[label=(\alph*),leftmargin=2.5em]
\item the causal relation $\leq$ on the set $M$ of points is a partial order (reflexive, transitive
and antisymmetric);
\item the support map $h\mapsto\supp h$ from $\cF$ to subsets of $M$ satisfies $\supp0=\emptyset$,
$\supp(h+h')\subseteq\supp h\cup\supp h'$ and $\supp(\lambda h)\subseteq\supp h$ for $\lambda\in\R$, and
each $\supp h$ is a closed subset of $M$;
\item the wave operator $P$ maps $\cF$ into $\cF$ and $\cD$ into $\cD$, and $P(\iota f)=\iota(Pf)$ for
$f\in\cD$ (so $P_\cD$ is $P$ on test functions);
\item the pairing $\langle f,h\rangle=\int_Mfh\,dp$ of $f\in\cD$ and $h\in\cF$ vanishes when
$\supp f\cap\supp h=\emptyset$, and it is nondegenerate in $h$: if $\langle f,h\rangle=0$ for all
$f\in\cD$, then $h=0$;
\item (integration by parts) for $h,h'\in\cF$ whose supports meet in a compact subset of $M$,
$\int_M(Ph)h'\,dp=\int_Mh\,(Ph')\,dp$; in particular $\langle Pf,h\rangle=\langle f,Ph\rangle$ for
$f\in\cD$ and $h\in\cF$.
\end{enumerate}
\end{proposition}
\begin{proof}
\begin{enumerate}[leftmargin=2.5em]
\item[(a)] Reflexivity is $0\leq0$; transitivity is the triangle inequality; if $p\leq q\leq p$ then
$q^0=p^0$ and $|\vec q-\vec p|\le0$.
\item[(b)] $\{h+h'\neq0\}\subseteq\{h\neq0\}\cup\{h'\neq0\}$; take closures.
\item[(c)] Derivatives of smooth functions are smooth and do not enlarge supports.
\item[(d)] Locality: the integrand vanishes everywhere. Nondegeneracy: if $\int fh=0$ for all
$f\in\cD$, then $h=0$ almost everywhere by the fundamental lemma of the calculus of variations, hence
everywhere by continuity.
\item[(e)] Integrating by parts twice along each coordinate direction gives the identity when one of
$h$, $h'$ has compact support. In general both integrands vanish outside the compact set
$\supp h\cap\supp h'$; multiplying $h'$ by a test function equal to $1$ near that set changes neither
integrand and makes $h'$ compactly supported.
\end{enumerate}
\end{proof}

\subsection{The retarded and advanced Green operators}

\begin{definition}[Cone measures and Green operators]\label{def:kirchhoff}
The \emph{future cone measure} $\mu^+$ on $M$ is the image of the measure $dy/(4\pi|y|)$ on $\R^3$
under $y\mapsto(|y|,y)$, and the \emph{past cone measure} $\mu^-$ is the image of $\mu^+$ under
$w\mapsto-w$. For $f\in\cD$,
\[
(G^\pm f)(p)=\int_Mf(p-w)\,\mu^\pm(dw).
\]
\end{definition}

\begin{proposition}[\revm{The retarded potential}]\label{prop:kirchhoff}
For $f,g\in\cD\revk{=C_c^\infty(M)}$, \revk{the Green operators $G^\pm$ of Definition~\ref{def:kirchhoff}
satisfy}:
\begin{enumerate}[label=(\alph*),leftmargin=2.5em]
\item $(G^\pm f)(t,x)=\frac1{4\pi}\int_{\R^3}f(t\mp|y|,x-y)\,|y|^{-1}\,dy$;
\item $G^\pm f$ is smooth, $G^\pm:\cD\to\cF$ is linear, and $\partial_vG^\pm f=G^\pm\partial_vf$ for
every direction $v$; hence $PG^\pm f=G^\pm Pf$;
\item $\supp G^\pm f\subseteq J^\pm(\supp f)$;
\item $\langle f,G^+g\rangle=\langle g,G^-f\rangle$.
\end{enumerate}
\end{proposition}
\begin{proof}
\begin{enumerate}[leftmargin=2.5em]
\item[(a)] Integration against an image measure; for $G^-$ it gives $f(t+|y|,x+y)$, and the
substitution $y\mapsto-y$ gives the formula.
\item[(b)] $\mu^\pm$ are finite on bounded sets, since $\int_{|y|<r}dy/(4\pi|y|)=r^2/2$. So $G^\pm f$ is
the convolution of the smooth, compactly supported $f$ with a locally finite measure, which is
smooth. The derivative is taken under the integral, dominated by a constant on a compact set.
\item[(c)] If $(G^+f)(t,x)\neq0$, the integrand of (a) is not identically zero, so
$(t-|y|,x-y)\in\supp f$ for some $y$, a point below $(t,x)$ on its light cone. So
$\{G^+f\neq0\}\subseteq J^+(\supp f)$, which is closed by Lemma~\ref{lem:mink-compact}(b); likewise
for $G^-$.
\item[(d)] $(p,w)\mapsto f(p)g(p-w)$ is continuous and vanishes unless $p\in\supp f$ and
$w\in p-\supp g$, a compact set on which $dp\,\mu^+(dw)$ is finite. So Fubini applies, and with
translation invariance of Lebesgue measure both sides equal $\iint f(p)g(p-w)\,dp\,\mu^+(dw)$.
\end{enumerate}
\end{proof}
\noindent\motiv{$G^+$ samples $f$ in the causal past of the point and is supported in the causal
future of $\supp f$: it is the retarded operator in the naming of Remark~\ref{rem:sign-table}.}

\subsection{The fundamental solution}

\begin{theorem}[\revm{The fundamental solution}]\label{thm:mink-core}
For \revk{every test function} $f\in\revk{C_c^\infty(M)}$ \revk{and the wave operator
$P=\partial_t^2-\nabla^2$},
\[
\int_{\R^3}\frac{(Pf)(|y|,y)}{|y|}\,dy=4\pi\,f(0).
\]
\end{theorem}
\begin{proof}
\emph{Spherical means.} Let $(\mathsf Af)(s,y)$ be the mean of $f(s,\cdot)$ over the sphere of
radius $|y|$ about $0$. The weight $1/|y|$ and Lebesgue measure are rotation invariant, so the left
side does not change when $Pf$ is replaced by $\mathsf APf$; and $\mathsf APf=P\mathsf Af$, because
rotations commute with $\nabla^2$. Put $\bar f(s,\rho)=(\mathsf Af)(s,\rho e_1)$ for $\rho\in\R$, with
$e_1$ a unit vector of $\R^3$. It is smooth and compactly supported on $\R^2$, $(\mathsf
Af)(s,y)=\bar f(s,|y|)$, and $\bar f(0,0)=f(0)$.

\emph{Radial calculus.} For $y\neq0$, $P[\bar f(s,|y|)]=(\bar f_{ss}-\bar f_{\rho\rho}-\tfrac2\rho
\bar f_\rho)(s,|y|)$ with $\rho=|y|$, and $\int_{\R^3}\Phi(|y|)\,dy=4\pi\int_0^\infty\rho^2\Phi(\rho)\,d\rho$
for integrable $\Phi(|y|)$. The left side becomes
$4\pi\int_0^\infty\rho\,(\bar f_{ss}-\bar f_{\rho\rho}-\tfrac2\rho\bar f_\rho)(\rho,\rho)\,d\rho$.

\emph{A null ray.} $\rho(\bar f_{ss}-\bar f_{\rho\rho}-\tfrac2\rho\bar f_\rho)=(\partial_s+
\partial_\rho)(\partial_s-\partial_\rho)(\rho\bar f)$, and $d/d\rho$ along the diagonal $s=\rho$ is
$\partial_s+\partial_\rho$. So the integrand is the derivative along the diagonal of
$(\partial_s-\partial_\rho)(\rho\bar f)$, which vanishes for large $\rho$ and equals $-\bar f(0,0)$ at
$\rho=0$. The integral is $\bar f(0,0)=f(0)$.
\end{proof}

\begin{remark}[The method of spherical means]\label{rem:spherical-means}
The theorem says that $\delta(t-|y|)/(4\pi|y|)$ is a fundamental solution of $P$. The proof is the
classical method of spherical means~\cite[\S2.4.1]{evans}, which reduces the wave equation in $3+1$
dimensions to the one-dimensional wave equation for $\rho\bar f$.

\artefact{The formalization writes the spherical mean as an average over the group $O(3)$ of
orthogonal operators on $\R^3$ with its Haar probability measure, $(\mathsf
Af)(s,y)=\int_{O(3)}f(s,Q^{-1}y)\,dQ$. Mathlib has Haar measure on compact groups; for its measure on
spheres it lacks the rotation invariance, and it lacks the divergence theorem on balls\revm{, on which
the textbook proof rests}.} $O(3)$ acts transitively on each sphere: if
$|y|=|y'|$ and $y\neq y'$, the reflection \revk{$1-2zz^\top/|z|^2$ with $z=y'-y$} maps $y'$ to $y$.
\end{remark}

\begin{corollary}[The Green identities]\label{cor:mink-green}
For $f\in\cD$, \revk{the operators $G^\pm$ of Definition~\ref{def:kirchhoff} and the wave
operator $P$ satisfy} $G^+Pf=f$, $G^-Pf=f$, $PG^+f=f$ and $PG^-f=f$.
\end{corollary}
\begin{proof}
$P$ commutes with translations and with the reflections of time and of space. Fix $(t,x)$ and let
$f_{t,x}(s,z)=f(t-s,x-z)$; then $(Pf_{t,x})(|y|,y)=(Pf)(t-|y|,x-y)$. So by
Proposition~\ref{prop:kirchhoff}(a) and Theorem~\ref{thm:mink-core},
$(G^+Pf)(t,x)=f_{t,x}(0)=f(t,x)$. For $G^-$ use $(s,z)\mapsto f(t+s,x-z)$. Finally,
$PG^\pm f=G^\pm Pf$ by Proposition~\ref{prop:kirchhoff}(b).
\end{proof}

\subsection{Cutoffs and the Green data}

\noindent\motiv{\emph{What the cutoffs do.} The observable $W([f])$ depends on the test function
$f$ only through the solution $Ef$ it generates, since $[f]$ is the class of $f$ modulo $\ker E$. The
time-slice axiom asks that every such observable already belong to the algebra of a thin slab of
time $S$. To show this, one moves the smearing of $f$ into $S$ without changing $Ef$.}

\medskip\noindent\motiv{Take $\chi$ equal to $0$ before $S$ and to $1$ after it, and replace $f$ by
$f'=-P(\chi Ef)$. Since $PEf=0$, the function $P(\chi Ef)$ can be nonzero only where $\chi$ changes,
that is inside $S$, and only where $Ef$ lives, inside the causal shadow of $\supp f$: this is
condition (c1). And $f-f'=P\bigl((1-\chi)G^+f+\chi G^-f\bigr)$ is $P$ of a test function, which $E$
annihilates, so $Ef'=Ef$: this is condition (c2). Physically, $f'$ is $f$ transported along the
equation of motion into the slab.}

\medskip\noindent\motiv{The slab must be allowed to be thin and to sit anywhere, and the family of slabs must be
carried to itself by the symmetries. Here the slabs are the time intervals $t_1\le\tau\le t_2$ of all
inertial observers.}

\begin{definition}[Inertial slabs]\label{def:mink-cutoff}
A vector $n$ is \emph{future unit}\revk{, short for future-directed unit timelike,} if $n^0>0$ and $\eta(n,n)=-1$; its time function is
$\tau_n(p)=-\eta(n,p)=n^0p^0-\vec n\cdot\vec p$, the time of the inertial observer with four-velocity
$n$.

A \emph{cutoff} is $(n,t_1,t_2)$ with $n$ future unit and $t_1<t_2$. Its \emph{slab} is
$S=\{t_1\leq\tau_n\leq t_2\}$, and $\chi$ is multiplication by
$\chi(p)=\vartheta\bigl((\tau_n(p)-t_1)/(t_2-t_1)\bigr)$, with $\vartheta$ a smooth function equal to $0$
on $(-\infty,0]$ and to $1$ on $[1,\infty)$ (\artefact{Mathlib's smooth transition function}).
\end{definition}

\begin{lemma}[Time functions and compactness]\label{lem:mink-compact}
Let $n$ be future unit, \revk{$\tau_n$ its time function (Definition~\ref{def:mink-cutoff}), and
$J^\pm$ the causal future and past for the order $\leq$ of Definition~\ref{def:mink}}.
\begin{enumerate}[label=(\alph*),leftmargin=2.5em]
\item $n^0>|\vec n|$, and $\tau_n(v)\geq(n^0-|\vec n|)v^0\geq0$ for $v$ future causal.
\item For $Y\subseteq M$ compact, $J^\pm(Y)$ are closed, and $J^+(Y)\cap\{\tau_n\leq t_2\}$ and
$J^-(Y)\cap\{\tau_n\geq t_1\}$ are compact for all $t_1,t_2\in\R$.
\end{enumerate}
\end{lemma}
\begin{proof}
\begin{enumerate}[leftmargin=2.5em]
\item[(a)] $(n^0)^2=1+|\vec n|^2$, and $\tau_n(v)\geq n^0v^0-|\vec n||\vec v|\geq(n^0-|\vec n|)v^0$ by
Cauchy--Schwarz.
\item[(b)] $J^+(Y)$ is the sum of $Y$ and the closed cone of future causal vectors, which is closed.
For $Y\neq\emptyset$ and $q=y+v$ with $y\in Y$ and $\tau_n(q)\leq t_2$, (a) bounds $v^0\geq|\vec v|$ by
$(t_2-\min_Y\tau_n)/(n^0-|\vec n|)$, so the set is bounded; for $Y=\emptyset$ it is empty. The past
case follows by $p\mapsto-p$.
\end{enumerate}
\end{proof}

\begin{proposition}[Conditions (c1) and (c2)]\label{prop:mink-cutoff}
For every cutoff \revk{$(n,t_1,t_2)$ of Definition~\ref{def:mink-cutoff}, with slab $S$ and
multiplication $\chi$,} and every $f\in\cD$, with $E=G^--G^+$ \revk{for the operators of
Definition~\ref{def:kirchhoff}}:
\begin{enumerate}[label=(c\arabic*),leftmargin=2.5em]
\item $\supp P(\chi Ef)\subseteq S\cap J(\supp f)$;
\item $(1-\chi)G^+f+\chi G^-f\in\cD$.
\end{enumerate}
\end{proposition}
\begin{proof}
\emph{(c1).} $P$ is local: functions equal near $p$ have equal $P$ at $p$. Below the slab
$\chi Ef=0$ near $p$; above it $\chi Ef=Ef$ near $p$, and $PEf=f-f=0$ (Corollary~\ref{cor:mink-green}).
Outside $J(\supp f)$, the open cone $p+I^-$ misses $J^+(\supp f)$, since $p'\leq q$ and $q-p\in I^-$
would give $p'\leq p$ by transitivity. So $\chi G^+f=0$ there and $P(\chi G^+f)=0$ on $p+I^-$. The
points $p-se_0$, $s>0$, tend to $p$, so $P(\chi G^+f)(p)=0$ by continuity. Likewise
$P(\chi G^-f)(p)=0$ with the open cone $p+I^+$. So $P(\chi Ef)$ vanishes outside $S\cap J(\supp f)$, which is closed ($\tau_n$ is continuous;
Lemma~\ref{lem:mink-compact}(b)) and therefore contains the support.

\emph{(c2).} The function vanishes outside
$\bigl(J^+(\supp f)\cap\{\tau_n\leq t_2\}\bigr)\cup\bigl(J^-(\supp f)\cap\{\tau_n\geq t_1\}\bigr)$,
which is compact by Lemma~\ref{lem:mink-compact}.
\end{proof}

\begin{theorem}[The free massless field]\label{thm:mink-green}
\revm{The objects of} Definition~\ref{def:mink}, the operators $G^\pm$ \revk{of Definition~\ref{def:kirchhoff}} and the cutoffs of Definition~\ref{def:mink-cutoff}
\revm{form} Green data (Definition~\ref{def:green-data}). If the CCR algebra of the derived phase space has
an irreducible representation, the CCR net of the free massless field is generating, causal and
has the time-slice property over the derived causal index, and its algebra is primitive.
\end{theorem}
\begin{proof}
(G1)--(G5) are Proposition~\ref{prop:mink-basic}; (G6) is Proposition~\ref{prop:kirchhoff}(b)--(d)
and Corollary~\ref{cor:mink-green}; (G7) is Proposition~\ref{prop:mink-cutoff}. Then
Theorem~\ref{thm:green-axioms}.
\end{proof}
\noindent\motiv{As after Theorem~\ref{thm:green-axioms}, generation, causality and the time-slice
property hold without the hypothesis; only primitivity uses it. The data are not trivial: for $f\geq0$ nonzero, $G^+f$ vanishes before $\supp f$ while $G^-f$
does not, so $Ef\neq0$, $V\neq0$, and $\sigma\neq0$ by nondegeneracy
(Proposition~\ref{prop:green-nondeg}). Regions are all subsets of $M$. One whose interior is empty contains the support of no nonzero
test function, so its algebra is $\C1$.}

\begin{remark}[Which time slices]\label{rem:mink-slabs}
The cutoffs are slabs between parallel spacelike hyperplanes, one family
for each inertial frame $n$. Remark~\ref{rem:green-continuum} uses all Cauchy temporal functions, of
which each $\tau_n$ is one. Fewer cutoffs give fewer pairs $O\dep O'$, so the derived dependence here
is contained in the one derived from all Cauchy temporal functions (both with closed supports).
\begin{itemize}[leftmargin=2em]
\item \emph{Flat Cauchy surfaces.} For $\Sigma=\{\tau_n=c\}$, the first argument of
Remark~\ref{rem:green-dep-spacetime} runs with $\mathcal T=\tau_n-c$ and
Lemma~\ref{lem:mink-compact}(b). So\revk{, for open regions with compact closure,} the dependence
of Example~\ref{ex:gh-index} with such a $\Sigma$ implies the one here, and the time-slice axiom in
that form holds.
\item \emph{Curved Cauchy surfaces.} Here \revk{the implication} can fail.

\revk{\emph{Dependence in the sense of Example~\ref{ex:gh-index}.}} Let $\phi$ be smooth on $\R^3$
with $\phi(\vec x)=|\vec x|^2$ for $|\vec x|\le2$ and $\revk{|\nabla\phi|\le5}$, so that
$\Sigma=\{t=\phi(\vec x)/10\}$ is a smooth spacelike Cauchy surface. \revk{Such $\phi$ exist: take
$\phi(\vec x)=g(|\vec x|^2)$ with $g$ smooth, $g(s)=s$ for $s\le4$, and $g'$ decreasing from $1$ to $0$
on $[4,5]$ and zero beyond; then $|\nabla\phi|=2|\vec x|\,g'(|\vec x|^2)\le2\sqrt5$.} Let
$D=\{|t|+|\vec x|<1\}$ and $O'=\{|t-\phi(\vec x)/10|<1/40,\ |\vec x|<2\}$, both open with compact
closure.

A causal curve from $D$ meets $\Sigma$ where $|\vec x|<1+\revk{|\phi(\vec x)|}/10$\revk{. For
$|\vec x|\ge2$, $|\phi(\vec x)|\le4+5(|\vec x|-2)$ would give $|\vec x|<4/5$, which is impossible; so
$|\vec x|<2$, $\phi(\vec x)=|\vec x|^2$, and} $|\vec x|<1.13$. So every endless causal curve through $D$ meets
$\Sigma$ in $O'$, and $D\dep O'$ in the sense of Example~\ref{ex:gh-index}.

\revk{\emph{No inertial slab serves.}} Now take $f\in\cD(D)$ with
$\{0\}\times\{|\vec x|\le3/4\}\subseteq\supp f$. Then $J(\supp f)$ contains every line
$\{(t,\vec x):t\in\R\}$ with $|\vec x|\le3/4$, and a slab meets such a line in a segment whose lowest
point has $t$ an affine function of $\vec x$. For $S\cap J(\supp f)\subseteq O'$ that affine function
would have to stay within $1/40$ of $|\vec x|^2/10$ on the ball $|\vec x|\le3/4$. But every affine
function differs from $|\vec x|^2/10$ by at least $(3/4)^2/20=9/320>1/40$ somewhere on it (compare
$\vec x=0$ with $\vec x=\pm\tfrac34e_1$). So $D\dep O'$ fails \revk{for the inertial slabs}.

\revk{The inclusion $\cA(D)\subseteq\cA(O')$ itself still holds, on paper. The local algebras do
not depend on the cutoffs. With all Cauchy temporal functions as cutoffs,
Remark~\ref{rem:green-dep-spacetime} gives $D\dep O'$ for the derived index,
Proposition~\ref{prop:green-phase} gives phase-space data over it, and the time-slice property of
Theorem~\ref{thm:weyl-net} gives the inclusion. It is not a consequence of
Theorem~\ref{thm:mink-green}.}
\end{itemize}
Cutoffs from curved Cauchy temporal functions \revm{are not treated in these notes}. The family of inertial slabs is
carried to itself by Poincar\'e maps (Proposition~\ref{prop:mink-symmetry}), which covariance needs.
\revk{Neither bullet is formalized.}
\end{remark}

\subsection{Uniqueness and Poincar\'e covariance}

\begin{theorem}[Uniqueness of the Green operators]\label{thm:mink-unique}
If \revk{a map} $G':\cD\to\cF$ satisfies $PG'f=f$ and $\supp G'f\subseteq J^+(\supp f)$ for every
$f\revk{\in\cD}$, then $G'$ \revk{is the operator $G^+$ of Definition~\ref{def:kirchhoff}}; likewise with
$J^-$ and $G^-$.
\end{theorem}
\begin{proof}
\revk{\emph{Compact overlap.}} Let $f,g\in\cD$. The supports of $G^-g$ and $G'f$ lie in the closed sets $J^-(\supp g)$ and
$J^+(\supp f)$. If \revk{$q\leq q'$ then $q^0\leq q'^0$}, so $J^-(\supp g)\subseteq\{\tau_{e_0}\leq t_2\}$ with $t_2$
\revm{an upper bound of the times in} $\supp g$. The intersection is therefore a closed subset of
$J^+(\supp f)\cap\{\tau_{e_0}\leq t_2\}$, which is compact by Lemma~\ref{lem:mink-compact}(b).

\revk{\emph{Pairing.}} By Proposition~\ref{prop:mink-basic}(e), Corollary~\ref{cor:mink-green} and
Proposition~\ref{prop:kirchhoff}(d),
\[
\langle g,G'f\rangle=\int(PG^-g)G'f=\int G^-g\,PG'f=\langle f,G^-g\rangle=\langle g,G^+f\rangle,
\]
and nondegeneracy gives $G'f=G^+f$.
\end{proof}

\begin{definition}[Orthochronous Poincar\'e group]\label{def:mink-poincare}
An \emph{orthochronous Lorentz map} is a linear $\Lambda:\R^4\to\R^4$ with
$\eta(\Lambda u,\Lambda w)=\eta(u,w)$ and $(\Lambda e_0)^0>0$. The \emph{orthochronous Poincar\'e
group} $\mathcal P^\uparrow$ consists of the maps $\gamma(p)=\Lambda p+a$, $a\in\R^4$, under composition;
it contains the space reflections.
\end{definition}

\begin{lemma}[Orthochronous Lorentz maps]\label{lem:mink-lorentz}
Let $\gamma(p)=\Lambda p+a$ be in $\mathcal P^\uparrow$.
\begin{enumerate}[label=(\alph*),leftmargin=2.5em]
\item $\Lambda$ is invertible, $\Lambda^{-1}$ and composites of orthochronous Lorentz maps are
orthochronous Lorentz maps, and $\mathcal P^\uparrow$ is a group.
\item $\gamma p\leq\gamma q$ iff $p\leq q$.
\item If $n$ is future unit, so is $\Lambda n$, and \revk{$\tau_n(\Lambda^{-1}p)=\tau_{\Lambda n}(p)$}.
\item $|\det\Lambda|=1$, so $\gamma$ preserves Lebesgue measure.
\item $P(h\circ\gamma)=(Ph)\circ\gamma$ for smooth $h$.
\end{enumerate}
\end{lemma}
\begin{proof}
\begin{enumerate}[leftmargin=2.5em]
\item[(a)] $\eta$ is nondegenerate, so $\Lambda u=0$ implies $u=0$; $\Lambda^{-1}$ and composites
preserve $\eta$. Orthochrony of $\Lambda^{-1}$ and of composites is shown in (b) and (c).
\item[(b)] Let $n'=\Lambda^{-1}e_0$. Then $\eta(n',n')=\eta(e_0,e_0)=-1$ and
$n'^0=-\eta(\Lambda^{-1}e_0,e_0)=-\eta(e_0,\Lambda e_0)=(\Lambda e_0)^0>0$, so $n'$ is future unit and
$\Lambda^{-1}$ is orthochronous. For $v$ future causal,
$(\Lambda v)^0=-\eta(\Lambda v,e_0)=-\eta(v,n')=\tau_{n'}(v)\geq(n'^0-|\vec n'|)v^0\geq0$ by
Lemma~\ref{lem:mink-compact}(a), and $\eta(\Lambda v,\Lambda v)=\eta(v,v)\leq0$; so $\Lambda v$ is future
causal. The converse is the same for $\Lambda^{-1}$.
\item[(c)] $\eta(\Lambda n,\Lambda n)=-1$, and $(\Lambda n)^0\geq(n'^0-|\vec n'|)n^0>0$ by the inequality
in (b), since $n^0>|\vec n|$. With $n=\Lambda'e_0$ this gives orthochrony of $\Lambda\Lambda'$. Also
\revk{$-\eta(n,\Lambda^{-1}p)=-\eta(\Lambda n,p)$}.
\item[(d)] In the standard basis $\Lambda^\top\eta\Lambda=\eta$ and $\det\eta=-1$, so
$(\det\Lambda)^2=1$.
\item[(e)] $Ph=-\operatorname{tr}(\eta^{-1}\partial^2h)$ with $\partial^2h$ the matrix of second
derivatives, which transforms as $\Lambda^\top\partial^2h\,\Lambda$. And
$\Lambda\eta^{-1}\Lambda^\top=\eta^{-1}$ follows from $\Lambda^\top\eta\Lambda=\eta$: here $\eta^{-1}=\eta$,
$\Lambda^{-1}=\eta\Lambda^\top\eta$, so $\Lambda\eta\Lambda^\top\eta=1$. Translations commute with $P$.
\end{enumerate}
\end{proof}

\begin{proposition}[Poincar\'e maps are symmetries]\label{prop:mink-symmetry}
Every $\gamma=(\Lambda,a)\in\mathcal P^\uparrow$ gives a symmetry of the Green data
(Definition~\ref{def:green-symmetry}) with point map $\gamma$, $\gamma_\cD f=f\circ\gamma^{-1}$,
$\gamma_\cF h=h\circ\gamma^{-1}$, and on cutoffs
$(n,t_1,t_2)\mapsto(\Lambda n,\,t_1+\tau_{\Lambda n}(a),\,t_2+\tau_{\Lambda n}(a))$.
\end{proposition}
\begin{proof}
\emph{Order, supports and pairing.} The order is preserved by Lemma~\ref{lem:mink-lorentz}(b);
$\gamma$ is a homeomorphism, so it maps supports to supports; and the pairing is invariant by (d).

\emph{Slabs.} By (c), $\gamma(S)=\{q:t_1\le\tau_n(\Lambda^{-1}(q-a))\le t_2\}$ is the slab of the image
cutoff. The cutoff map is a bijection, with inverse the map of $\gamma^{-1}$, because the shifts add
under composition: $\tau_{\Lambda\Lambda'n}(\Lambda a'+a)=\tau_{\Lambda'n}(a')+\tau_{\Lambda\Lambda'n}(a)$ by
(c).

\emph{Green operators.} $G'=\gamma_\cF G^+\gamma_\cD^{-1}$ satisfies $PG'f=f$ by (e) and
$\supp G'f\subseteq J^+(\supp f)$ by (b), so $G'=G^+$ by Theorem~\ref{thm:mink-unique}; likewise for
$G^-$.
\end{proof}

\begin{theorem}[Poincar\'e covariance]\label{thm:mink-covariance}
The CCR net of Theorem~\ref{thm:mink-green} is covariant under $\mathcal P^\uparrow$: with
$\kappa_\gamma(O)=\gamma(O)$ and $\alpha_\gamma$ the automorphism of the CCR algebra induced by the
symmetry of Proposition~\ref{prop:mink-symmetry}, $\gamma\mapsto\kappa_\gamma$ and
$\gamma\mapsto\alpha_\gamma$ are group homomorphisms and every $(\kappa_\gamma,\alpha_\gamma)$ is a
covariance pair (Definition~\ref{def:covariance}).
\end{theorem}
\begin{proof}
Proposition~\ref{prop:green-symmetry} gives, for each $\gamma$, an isomorphism of phase-space data with
$L_\gamma[f]=[\gamma_\cD f]$. Since $\gamma_\cD$ composes as $\gamma$ does,
$L_{\gamma\gamma'}=L_\gamma L_{\gamma'}$ and $L_{\mathrm{id}}=\mathrm{id}$; likewise
$\kappa_{\gamma\gamma'}=\kappa_\gamma\kappa_{\gamma'}$. Theorem~\ref{thm:ccr-cov} turns these
isomorphisms into covariance pairs with $\alpha_{\mathrm{id}}=\mathrm{id}$ and
$\alpha_{\gamma\gamma'}=\alpha_\gamma\alpha_{\gamma'}$.
\end{proof}
\noindent\motiv{Every isometry of Minkowski space is affine, so $\mathcal P^\uparrow$ is the group of
isometries that preserve the time orientation. This is covariance in the sense of
Definition~\ref{def:covariance} under that group, Axiom~5 of~\cite{dimock} for this spacetime\revo{, restricted to the isometries that preserve the time orientation (Remark~\ref{rem:time-orientation})}.}

\medskip\noindent\motiv{Two other maps are left out. Time reversal exchanges $G^+$ and $G^-$, so it is not a symmetry
of Green data in the sense of Definition~\ref{def:green-symmetry}. Dilations $p\mapsto\lambda p$, $\lambda>0$, also preserve the order and the slabs, and, with
$f\mapsto\lambda^{-3}f(\cdot/\lambda)$ on test functions, the form $\sigma$. Definition~\ref{def:green-symmetry}
does not express rescalings, and they are not used.}
\subsection{Green operators of a local operator}\label{sec:kg-greenops}

\noindent\motiv{The proofs of \revm{this chapter} on cutoffs, uniqueness and covariance use
only these properties of the wave operator and its Green operators: the operator is local and can be
integrated by parts; its Green operators satisfy the Green identities, are formally adjoint and are
supported in the causal future and past; and, for covariance, the operator is Poincar\'e invariant.}

\medskip\noindent\motiv{This section states the construction on Minkowski space for every operator with these properties
(Theorem~\ref{thm:kg-greenops}), so that for the Klein--Gordon operator only the Green operators
remain to be constructed. In this section $P$ is any such operator and $G^\pm$ are its Green
operators\revm{; in Chapter~\ref{sec:mass}}, from Definition~\ref{def:kg-supported} on, $P$ is the wave operator and $G^\pm$ are the
extended Green operators of Definition~\ref{def:kg-supported}. \artefact{In the formalization the results of
\revm{this chapter} on cutoffs, uniqueness and covariance are derived as the instance of
Theorem~\ref{thm:kg-greenops} in Proposition~\ref{prop:kg-wave}, with their statements
unchanged.}}

\begin{definition}[Local operators and their Green operators]\label{def:kg-greenops}
Let $M$, $\le$, $\cD$, $\cF$, $\iota$, $\supp$ and the pairing be as in Definition~\ref{def:mink}. A
linear map $P:\cF\to\cF$ is \emph{local} if, for all $h,h'\in\cF$ and $p\in M$, $h=h'$ on a
neighbourhood of $p$ implies $(Ph)(p)=(Ph')(p)$. It is \emph{Poincar\'e invariant} if
$P(h\circ\gamma^{-1})=(Ph)\circ\gamma^{-1}$ for all $h\in\cF$ and all $\gamma\in\mathcal P^\uparrow$
(Definition~\ref{def:mink-poincare}).

Let $P$ be local, and let $P_\cD:\cD\to\cD$ be linear with $\iota P_\cD=P\iota$ and
$\langle P_\cD f,h\rangle=\langle f,Ph\rangle$ for $f\in\cD$ and $h\in\cF$ (formal self-adjointness).
Since $\iota$ is injective, $P_\cD$ is determined by $P$. \emph{Green operators for $P$} are linear
maps $G^\pm:\cD\to\cF$ such that, for all $f,g\in\cD$:
\begin{enumerate}[label=(\alph*),leftmargin=2.5em]
\item $PG^\pm f=f$ and $G^\pm P_\cD f=f$;
\item $\supp G^+f\subseteq J^+(\supp f)$ and $\supp G^-f\subseteq J^-(\supp f)$;
\item $\langle f,G^+g\rangle=\langle g,G^-f\rangle$.
\end{enumerate}
\revm{Condition (c) and $G^\pm P_\cD f=f$ follow from $PG^\pm f=f$, (b), locality and formal
self-adjointness, on paper: $\langle f,G^+g\rangle=\int(PG^-f)\,G^+g=\int G^-f\,PG^+g=\langle
g,G^-f\rangle$ by Lemma~\ref{lem:kg-ibp}, and then $\langle g,G^+P_\cD f\rangle=\langle P_\cD
f,G^-g\rangle=\langle f,PG^-g\rangle=\langle g,f\rangle$ for all $g$, so $G^+P_\cD f=f$ by
nondegeneracy (likewise for $G^-$). They are listed because Definition~\ref{def:green-data} lists
them, and for the Klein--Gordon operator they are proved directly.}
\end{definition}

\begin{lemma}[Locality and integration by parts]\label{lem:kg-ibp}
Let $P:\cF\to\cF$ be local (Definition~\ref{def:kg-greenops}). Then $\supp Ph\subseteq\supp h$ for
every $h\in\cF$. For $h_1,h_2\in\cF$ whose supports meet in a compact set, $(Ph_1)h_2$ and $h_1\,Ph_2$
are integrable on $M$; if moreover $P_\cD$ is as in Definition~\ref{def:kg-greenops}, then
$\int_M(Ph_1)h_2\,dp=\int_Mh_1\,(Ph_2)\,dp$.
\end{lemma}
\begin{proof}
\emph{Supports.} If $p\notin\supp h$, then $h=0$ near $p$, so $(Ph)(p)=(P0)(p)=0$ by locality. The
complement of $\supp h$ is open, so $Ph$ vanishes on it, and $\supp Ph\subseteq\supp h$.

\emph{Integrability.} Let $K=\supp h_1\cap\supp h_2$. By the first part both products are continuous
and vanish outside $K$, which is compact.

\emph{Integration by parts.} Take $\psi\in\cD$ equal to $1$ on a neighbourhood $U$ of $K$. On $U$,
$\psi h_1=h_1$, so $P(\psi h_1)=Ph_1$ there by locality; off $K$ both $(Ph_1)h_2$ and
$(P(\psi h_1))h_2$ vanish. Hence $(Ph_1)h_2=(P(\psi h_1))h_2$, and likewise $h_1\,Ph_2=(\psi
h_1)\,Ph_2$. With $\iota P_\cD=P\iota$ and formal self-adjointness,
$\int(Ph_1)h_2=\langle P_\cD(\psi h_1),h_2\rangle=\langle\psi h_1,Ph_2\rangle=\int h_1\,Ph_2$.
\end{proof}

\begin{theorem}[Green data from Green operators]\label{thm:kg-greenops}
Let $G^\pm$ be Green operators for a local operator $P$ (Definition~\ref{def:kg-greenops}).
\begin{enumerate}[label=(\alph*),leftmargin=2.5em]
\item The data of Definition~\ref{def:mink} with $P$, $P_\cD$, $G^\pm$ and the inertial slabs of
Definition~\ref{def:mink-cutoff} are Green data (Definition~\ref{def:green-data}). If the CCR algebra
of the derived phase space has an irreducible representation, the CCR net is generating, causal and
has the time-slice property over the derived causal index, and its algebra is primitive; as after
Theorem~\ref{thm:green-axioms}, only primitivity uses the hypothesis.
\item If a map $G':\cD\to\cF$ satisfies $PG'f=f$ and $\supp G'f\subseteq J^+(\supp f)$ for every
$f\in\cD$, then $G'=G^+$; likewise with $J^-$ and $G^-$.
\item If $P$ is Poincar\'e invariant, every $\gamma\in\mathcal P^\uparrow$ gives a symmetry of these
Green data with the maps of Proposition~\ref{prop:mink-symmetry}, and the CCR net is covariant under
$\mathcal P^\uparrow$ in the sense of Theorem~\ref{thm:mink-covariance}.
\end{enumerate}
\end{theorem}
\begin{proof}
\begin{enumerate}[leftmargin=2.5em]
\item[(a)] (G1)--(G5) are Proposition~\ref{prop:mink-basic}(a), (b), (d) and $\iota P_\cD=P\iota$,
and (G6) is Definition~\ref{def:kg-greenops}(a)--(c). The proof of Proposition~\ref{prop:mink-cutoff}
uses of $P$ only locality, continuity of $P(\chi G^\pm f)$ and $PEf=f-f=0$, together with the
supports of $G^\pm$ (Definition~\ref{def:kg-greenops}(b)) and the closedness of $J(\supp f)$, so it
gives (c1); (c2) uses no property of $P$. The rest is Theorem~\ref{thm:green-axioms}.
\item[(b)] The proof of Theorem~\ref{thm:mink-unique}, with Lemma~\ref{lem:kg-ibp} for the
integration by parts, Definition~\ref{def:kg-greenops}(a)--(c) for $G^\pm$, and
Proposition~\ref{prop:mink-basic}(d) for nondegeneracy.
\item[(c)] In the proof of Proposition~\ref{prop:mink-symmetry} only the Green operators involve $P$.
For them, the map $G'f=G^+(f\circ\gamma)\circ\gamma^{-1}$ of that proof satisfies $PG'f=f$ by
invariance and $\supp G'f\subseteq J^+(\supp f)$ by Lemma~\ref{lem:mink-lorentz}(b), so $G'=G^+$ by
(b); likewise for $G^-$. The proof of Theorem~\ref{thm:mink-covariance} does not involve $P$.
\end{enumerate}
\end{proof}

\begin{proposition}[The wave operator]\label{prop:kg-wave}
The wave operator $P$ of Definition~\ref{def:mink} is local and Poincar\'e invariant, and the
operators $G^\pm$ of Definition~\ref{def:kirchhoff} are Green operators for it
(Definition~\ref{def:kg-greenops}). Their Green data (Theorem~\ref{thm:kg-greenops}) are the Green
data of Theorem~\ref{thm:mink-green}.
\end{proposition}
\begin{proof}
\emph{Locality and invariance.} $(Ph)(p)$ is a combination of second derivatives at $p$, so $P$ is
local; it is Poincar\'e invariant by Lemma~\ref{lem:mink-lorentz}(e).

\emph{Green operators.} $\iota P=P\iota$ and formal self-adjointness are
Proposition~\ref{prop:mink-basic}(c), (e), and (a)--(c) of Definition~\ref{def:kg-greenops} are
Corollary~\ref{cor:mink-green} and Proposition~\ref{prop:kirchhoff}(c), (d). The two Green data
coincide component by component.
\end{proof}
\endgroup
